\begin{explanationbox}
\textbf{一、先把四个“心”看成四种不同的需求}
\begin{itemize}
  \item \textbf{重心 $G$：找平均位置。} 三个顶点的坐标各加起来除以 $3$。由于两个焦点 $(-c,0)$ 与 $(c,0)$ 正好抵消，重心就是 $P$ 的坐标缩小为三分之一。
  \item \textbf{内心 $I$：找离三条边一样远的点。} 内心是内切圆的圆心。这里最有价值的不是死记坐标，而是研究\textbf{内切圆在底边 $F_1F_2$ 上的切点为什么不动}。
  \item \textbf{垂心 $H$：找三条高的交点。} 底边 $F_1F_2$ 水平，因此从 $P$ 向底边作高必是竖线 $x=x_0$，垂心横坐标立刻得到。
  \item \textbf{外心 $O'$：找离三个顶点一样远的点。} 焦点关于 $y$ 轴对称，因此底边的垂直平分线是 $y$ 轴，外心横坐标立刻为 $0$。
\end{itemize}

\textbf{二、最直观的核心：为什么内心被“锁”在 $x=\pm a$？}

设内切圆与底边 $F_1F_2$ 切于 $T$。
内心 $I$ 到水平底边的垂线就是 $IT$，因此 $x_I=x_T$。
再利用\textbf{从同一个顶点向圆引出的两条切线长度相等}，可得
\[
F_1T=\frac{F_1F_2+PF_1-PF_2}{2},\qquad
x_T=-c+F_1T=\frac{PF_1-PF_2}{2}.
\]
右支上，双曲线第一定义告诉我们 $PF_1-PF_2=2a$，因此 $x_I=a$；
左支上两距离交换，故 $x_I=-a$。

\textbf{一句话理解：}\textbf{双曲线把“到两个焦点的距离差”锁死了，内切圆在焦点连线上的切点也就被锁死了。}
内心只能沿切点上方或下方的竖线运动，而不能左右移动。

\textbf{三、为什么纵坐标反而会变？}

设三角形面积为 $S$、半周长为 $s$、内切圆半径为 $r$。
由“将三角形分成三个以 $r$ 为高的小三角形”，有 $S=rs$。
因为底边长 $2c$、高为 $|y_0|$，所以 $S=c|y_0|$；
而双曲线的两条焦半径之和为 $2e|x_0|$，所以 $s=c+e|x_0|$。
因此
\[
r=\frac{c|y_0|}{c+e|x_0|},\qquad
 y_I=\operatorname{sgn}(y_0)r
     =\frac{ay_0}{a+|x_0|}.
\]
内心在 $x$ 方向不动，但会随三角形高度变化而上下运动。

\textbf{四、四心的位置，可以先不计算就猜出一半}
\begin{itemize}
  \item 焦点互为关于原点的对称点：$G_x=x_0/3$。
  \item 底边水平：$H_x=x_0$。
  \item 底边中垂线为 $y$ 轴：$O'_x=0$。
  \item 两焦半径之差固定：$I_x=\pm a$。
\end{itemize}
\textbf{所以应优先记住横坐标的结构，再求纵坐标。}

\textbf{五、与椭圆的对比：同一个几何公式，不同的曲线定义}

对于任意焦点三角形都成立的结论是
$\displaystyle x_I=(PF_1-PF_2)/2$。
\begin{itemize}
  \item \textbf{双曲线}：右支 $PF_1-PF_2=2a$，故 $x_I=a$（左支为 $-a$）；这是\textbf{距离差恒定}产生的定直线。
  \item \textbf{椭圆}：焦半径 $PF_1=a+ex_0$、$PF_2=a-ex_0$，故 $x_I=ex_0$；这里 $x_I$ 随 $P$ 变化。
  \item 椭圆还有 $PF_1+PF_2=2a$，因而 $y_I=cy_0/(a+c)$；双曲线则是 $y_I=cy_0/(c+e|x_0|)$。
\end{itemize}
\textbf{区别的本质：} 两种圆锥曲线的焦点三角形使用的是\textbf{同一套三角形几何}；
只是椭圆固定“距离和”，双曲线固定“距离差”，所以固定下来的内心坐标不同。

\textbf{六、内心不是在整条顶点切线上运动}

由双曲线方程，令 $u=|x_0|>a$，则
\[
|y_0|=\frac ba\sqrt{u^2-a^2},\qquad
|y_I|=\frac{a|y_0|}{a+u}
=b\sqrt{\frac{u-a}{u+a}}.
\]
所以 $0<|y_I|<b$。当 $P$ 无限靠近顶点时，内心趋近 $(\pm a,0)$；
当 $P$ 沿一支走向无穷远时，$|y_I|$ 趋近 $b$。
\textbf{顶点位置和 $y=\pm b$ 的端点都不是实际内心位置。}

\textbf{七、三个心之间的隐藏关系：欧拉线}

重心、外心、垂心满足
\[
\boxed{H=3G-2O'},\qquad \overrightarrow{O'H}=3\overrightarrow{O'G}.
\]
意思是 $O',G,H$ 在同一条直线上，而且 $O'G:GH=1:2$。
因此只要算出了外心，\textbf{垂心甚至不用再求高线}。
这不是双曲线专属，而是任意非退化三角形的欧拉线关系。
\end{explanationbox}
